% two_natured_electron.tex — the substrate-decided two-natured electron (Vol 9 Ch.3a)
% Included from chapters/03a_device_circuit_models.tex, alongside node_2domain_nport.tex.
% Renders the native-engine output of:
%   src/scripts/vol_9_device/two_natured_electron_figure.py
% Every curve is ACTUAL on-main engine output (no hand-drawn array). CONSISTENCY-
% class: re-expresses passing engine output + the already-canon two-natured
% structure (A1 mass-"3" / T2 charge-"3", A1 _|_ T2; S-matrix eigenvalues
% {+1,-1,-1,-1}). Not a new test, not a new claim.
\begin{figure}[htbp]
\centering
\includegraphics[width=\linewidth]{two_natured_electron_native_engine.png}
\caption{The substrate-decided two-natured electron, rendered from the native
engine (\texttt{src/scripts/vol\_9\_device/two\_natured\_electron\_figure.py}; every
curve is ACTUAL on-main engine output --- no hand-drawn array). This is the
two-port of \S\ref{sec:vol9_node_2domain_nport} read at the electron: the
charge-``3'' (Cosserat $(2,3)$ micro-rotation, $T_2$) and the mass-``3'' (A1
dilatation) are \textbf{two orthogonal natures} ($A_1\perp T_2$; the K4 S-matrix
eigenvalues $\{+1,-1,-1,-1\}$, already canon in \S\ref{sec:vol9_node_2domain_nport}
and Ch.~\ref{ch:vol9_pin_port_configuration}), \emph{not} one nature holding the
other.
\textbf{Panel A (static charge, cool).} The $(2,3)$ winding read as the boundary
linking integer $\mathcal{Q}=\mathrm{Link}(\partial\Omega,F)=3\in\mathbb{Z}$
(\texttt{charge\_quantization.seed\_pq\_winding}$\,\to\,$\texttt{compute\_Q\_link});
$w_{\mathrm{tor}}=2,\,w_{\mathrm{pol}}=3$, self-link $Q_{\mathrm H}=p\,q=6$. The
integer is UNCHANGED under four continuous topology-preserving deformations
(\texttt{deform\_continuous}: smooth-noise, swirl, local-scale, warp) --- pure
topology, no work. The charge is $\alpha$-\emph{free}: the integer comes from
\texttt{compute\_Q\_link}, never from $Q=1/\alpha$ or $Q_\text{TANK}=137$
(the $Q=137$ slot stays EMPTY).
\textbf{Panel B (dynamical mass, warm).} The A1 dilatation cavity mode, the
\textbf{$\Gamma=-1$-boundary-confined} dilatation breather (the surviving
localizer), eigenfrequency $\omega_{\mathrm{fork\,b}}=2.84$ (cold-cage anchor
$\approx2.84$/$2.87$), $a_{1}$ core-fraction $0.995$, lossless
($\mathrm{Im}(\omega)=0$); the winding-OFF HALT control recovers the same mode
(it is not a winding-coupling artifact). It is \emph{not} an autonomous bulk
self-trapping breather (that DISPERSED, MODE-III) and \emph{not} a
coupled-Hamiltonian-pinned bound mode (the joint dynamical locus tested NEGATIVE).
\textbf{mass $=$ A1 is the ratified grade-ASSIGNMENT (PR\#260), not a driver
measurement} --- no driver discriminates an A1-mass from a $T_2$-mass.
\textbf{Panel C (real-space body, neutral).} The real-space body is the $0_1$
unknot ($\mathcal{Q}_{link}=0$, $Q_{\mathrm H}=0$;
\texttt{cosserat\_field\_3d.initialize\_electron\_unknot\_sector}), trivial in
contrast to the $(2,3)$ ansatz on the same field type ($\mathcal{Q}_{link}=2$):
the $(2,3)$ is a PHASE-SPACE Clifford-torus $(V_{\mathrm{inc}},V_{\mathrm{ref}})$
winding, NOT a real-space knot
(\kbleaf{ave-kb/vol1/dynamics/ch4-continuum-electrodynamics/master-equation.md}:20,
the two-``3''s disambiguation; \texttt{cosserat\_field\_3d.py}:1176--1178,
Layer-3 phase-space). \textbf{Panels A and C draw the engine's seeded-frame
parametric tracers} (the toroidal geometry the field is laid on); \textbf{only
Panel B plots the raw field array}. CONSISTENCY-class re-expression of passing
engine output, not a new claim.}
\label{fig:vol9_two_natured_electron}
\end{figure}
